% TOP_PROBLEM: {"id": 1972, "title": "Erdős Problem #170", "category_id": 2, "category": {"id": 2, "name": "combinatorics", "display_name": "Combinatorics", "description": "Counting problems, graph theory, discrete structures.", "slug": "combinatorics", "order_index": 2, "created_at": "2026-07-31T15:26:25.671Z"}, "status": "open"}
% FIRST_POSTED: null
% ENTRY_KIND: statement_only
% Imported from: unsolvedmath_additional_500.tex; UnsolvedMath ID 1972
% !TeX program = lualatex
% Compile twice: lualatex 1972.tex
% Self-contained: no external bibliography, images, or \input files.
% Numeric section labels and visible numbers are original UnsolvedMath IDs.
\documentclass[11pt,a4paper]{article}
\usepackage[margin=24mm,headheight=14pt]{geometry}
\usepackage{fontspec}
\setmainfont{Latin Modern Roman}
\setsansfont{Latin Modern Sans}
\setmonofont{Latin Modern Mono}
\usepackage{amsmath,amssymb,mathtools}
\usepackage{unicode-math}
\setmathfont{Latin Modern Math}
\usepackage{microtype}
\usepackage{enumitem,xurl,needspace}
\usepackage[dvipsnames]{xcolor}
\usepackage{hyperref}
\hypersetup{unicode=true,colorlinks=true,linkcolor=MidnightBlue,urlcolor=MidnightBlue,
 pdftitle={UnsolvedMath Problem 1972},pdfauthor={Source-annotated compilation},bookmarksnumbered=true}
\usepackage{bookmark,tocloft,titlesec,fancyhdr}
\setlength{\cftsecnumwidth}{4.2em}
\cftsetpnumwidth{2.5em}
\cftsetrmarg{4em}
\titleformat{\section}{\Large\bfseries\raggedright}{\thesection}{0.7em}{}
\pagestyle{fancy}\fancyhf{}
\fancyhead[L]{\small\sffamily Additional UnsolvedMath statements}
\fancyhead[R]{\small Original catalogue IDs}
\fancyfoot[C]{\thepage}
\setlength{\parindent}{0pt}
\setlength{\parskip}{5pt plus 1pt}
\setlength{\emergencystretch}{3em}
\setcounter{tocdepth}{1}\setcounter{secnumdepth}{1}
\setlist[itemize]{leftmargin=3em,itemsep=4pt,topsep=4pt,parsep=0pt}
\allowdisplaybreaks
\newcommand{\N}{\mathbb{N}}
\newcommand{\Z}{\mathbb{Z}}
\newcommand{\Q}{\mathbb{Q}}
\newcommand{\R}{\mathbb{R}}
\newcommand{\C}{\mathbb{C}}
\newcommand{\F}{\mathbb{F}}
\newcommand{\A}{\mathbb{A}}
\newcommand{\PP}{\mathbb{P}}
\newcommand{\E}{\mathbb{E}}
\newcommand{\eps}{\varepsilon}
\newcommand{\dd}{\,\mathrm d}
\newcommand{\sref}[2]{\hyperlink{source:#1:#2}{[S#2]}}
\newif\ifshowcatalogue
\showcataloguetrue
\newcommand{\cataloguescope}[1]{\ifshowcatalogue
 \paragraph{Statement recorded in the supplied source.}
 #1\fi}
\begin{document}
% UnsolvedMath ID 1972; source catalogue code EP-170

\setcounter{section}{1971}

\section{The Optimal Constant for Interval Difference Bases}\label{1972}

{\small\sffamily UnsolvedMath ID 1972 · Combinatorics}

\subsection{Definitions and mathematical statement}

\paragraph{Shared notation.}Unless a section says otherwise, $\N=\{1,2,\ldots\}$,
$\Z$ is the ring of integers, $\Q,\R,\C$ are the rational, real, and complex
fields, $[N]=\{1,\ldots,\lfloor N\rfloor\}$, and $\log$ is the natural
logarithm. For real functions of a parameter tending to infinity,
$f\ll g$ and $f=O(g)$ mean $|f|\leq Cg$ eventually for some fixed $C$;
$f\gg g$ reverses this comparison; $f=o(g)$ means $f/g\to0$;
$f\sim g$ means $f/g\to1$; and $f\asymp g$ means both order bounds.
A subscript on $O$ or $\ll$ specifies allowed constant dependence.
The expression $n^{o(1)}$ in an upper bound means growth smaller than
$n^\epsilon$ for each fixed $\epsilon>0$. Local definitions govern any
exception, finite-field convention, or use of the same symbol for another
object. An asymptotic claim is not automatically an effective algorithm.



A difference basis here is constrained to $\{0,\ldots,N\}$ and must realize each target difference by two elements of the same set. Repetitions are allowed for the zero difference. The displayed question includes existence and identification of the limit. \sref{1972}{1} \sref{1972}{2}

\paragraph{Statement recorded in the supplied source.}
Let $F(N)$ be the smallest possible size of $A\subset \{0,1,\ldots,N\}$ such that $\{0,1,\ldots,N\}\subset A-A$. Find the value of $ \lim_{N\to \infty}\frac{F(N)}{N^{1/2}}. $ \sref{1972}{1}

{\small\emph{Transcription note.} Only unambiguous escape/formatting damage and nonmathematical serialization debris were repaired in the displayed source statement.}

\subsection{Short English statement}

What is the limiting square-root-normalized minimum size of a set whose pairwise differences cover an entire integer interval?

\subsection{Sources}

{\small

\begin{itemize}

\item[\hypertarget{source:1972:1}{S1}] ULAM.AI, \emph{UnsolvedMath}, supplied \texttt{problems.json}, record 1972 (EP-170), statement and background. The embedded literature-review date is 2026-08-17; that is the archive’s review date, not a fresh certification. Dataset landing page: \url{https://huggingface.co/datasets/ulamai/UnsolvedMath}.

\item[\hypertarget{source:1972:2}{S2}] T. F. Bloom, \emph{Erdős Problems}, Problem 170, \url{https://www.erdosproblems.com/170}. Reference imported from the supplied record; the live problem page was not independently checked for this compilation.

\item[\hypertarget{source:1972:3}{S3}] Erd\H{o}s, P. and G\'{a}l, I., On the representation of $1,2,\ldots,N$ by differences. Nederl. Akad. Wetensch., Proc. (1948), 1155-1158. Bibliographic information imported from the supplied record.

\item[\hypertarget{source:1972:4}{S4}] Leech, J., On the representation of $1,2,\ldots,n$ by differences. J. London Math. Soc. (1956), 160-169. Bibliographic information imported from the supplied record.

\item[\hypertarget{source:1972:5}{S5}] Pegg, E., Hitting All the Marks: Exploring New Bounds for Sparse Rulers and a Wolfram Language Proof. \url{https://blog.wolfram.com/2020/02/12/hitting-all-the-marks-exploring-new-bounds-for-sparse-rulers-and-a-wolfram-language-proof/} (2020). Bibliographic information imported from the supplied record.

\item[\hypertarget{source:1972:6}{S6}] Wichmann, B., A note on restricted difference bases. J. London Math. Soc. (1963), 465-466. Bibliographic information imported from the supplied record.

\end{itemize}

}

\label{end:1972}
\end{document}
